\documentclass[a4paper,10pt]{article}
\usepackage[czech]{babel}
\usepackage[T1]{fontenc}
\usepackage[utf8]{inputenc}
\usepackage{amsmath,amssymb}
\usepackage[margin=18mm]{geometry}
\usepackage{tikz}
\setlength{\parindent}{0pt}

\begin{document}
\section*{Komplexní čísla G/H -- řešení}
Podklad pro učitele k testu
\texttt{02\_test\_KomplexniCislaZaverecny\_GH.tex}.
Schválená doba je 45 minut, bez kalkulačky.

Zdrojové příklady jsou v
\texttt{ucebnice/06\_Komplexni\_cisla/Komplexní čísla.tex}:
G1 je úloha 45 g, G3 úloha 50 a, H3 úloha 50 e,
G5 úloha 53 d a H5 úloha 53 g.
H1 je vlastní obdobná úloha podle 45 g.
Obě zadání úlohy 2 jsou vlastní spojení geometrických podmínek
podle úlohy 23 e; obě zadání úlohy 4 jsou vlastní úlohy podle
převodů tvarů a Moivreovy věty (např. úloha 37).
Výběr, vlastní tvorbu a přiřazení schválil učitel v této konverzaci.

\subsection*{Varianta G}
\begin{enumerate}
\item Položme $z=x+yi$, kde $x,y\in\mathbb R$. Potom
$z\overline z=x^2+y^2$ a rovnice přejde na
\[
x^2+y^2-x-yi=6-2i.
\]
Porovnáním imaginárních částí dostáváme $y=2$;
pro reálné části pak $x^2-x-2=(x-2)(x+1)=0$.
Řešeními jsou $z=2+2i$ a $z=-1+2i$.
Dosazením ověříme
$|2+2i|^2-(2+2i)=8-2-2i=6-2i$ a
$|-1+2i|^2-(-1+2i)=5+1-2i=6-2i$.

\item Pro $z=x+yi$ je první podmínka rovností vzdáleností
od bodů $(1,1)$ a $(3,-1)$. Umocnění obou nezáporných stran dává
\[
(x-1)^2+(y-1)^2=(x-3)^2+(y+1)^2
\quad\Longleftrightarrow\quad x-y=2.
\]
Druhá podmínka je $x^2+y^2=10$. Dosadíme $x=y+2$:
\[
(y+2)^2+y^2=10
\quad\Longleftrightarrow\quad (y-1)(y+3)=0.
\]
Dostáváme $z=3+i$ a $z=-1-3i$. Oba body leží na kružnici
se středem v počátku a poloměrem $\sqrt{10}$ i na přímce $x-y=2$:
\begin{center}
\begin{tikzpicture}[scale=.58]
    \draw[step=1,gray!20,very thin] (-4,-4) grid (4,4);
    \draw[->] (-4.2,0)--(4.3,0) node[right] {$\operatorname{Re}z$};
    \draw[->] (0,-4.2)--(0,4.3) node[above] {$\operatorname{Im}z$};
    \draw (0,0) circle ({sqrt(10)});
    \draw[thick] (-1.5,-3.5)--(3.5,1.5) node[pos=.6,below right] {$x-y=2$};
    \fill (3,1) circle (2pt) node[above left] {$3+i$};
    \fill (-1,-3) circle (2pt) node[below right] {$-1-3i$};
\end{tikzpicture}
\end{center}

\item Diskriminant je
\[
D=(2-i)^2-4(3-i)=3-4i-12+4i=-9.
\]
Z $\sqrt D=3i$ vychází
\[
x_{1,2}=\frac{-2+i\pm3i}{2},
\qquad x_1=-1+2i,\quad x_2=-1-i.
\]
Kontrola Viètovými vztahy: součet je $-2+i$ a součin $3-i$.

\item Platí $1+\sqrt3i=2(\cos\frac\pi3+i\sin\frac\pi3)$ a
$1+i=\sqrt2(\cos\frac\pi4+i\sin\frac\pi4)$.
Podíl má modul $\sqrt2$ a argument $\frac\pi{12}$, proto
\[
\left(\frac{1+\sqrt3i}{1+i}\right)^6
=(\sqrt2)^6\left(\cos\frac\pi2+i\sin\frac\pi2\right)=8i.
\]
Nezávisle algebraicky: druhá mocnina podílu je $\sqrt3+i$,
jehož třetí mocnina je $8i$.

\newpage
\item Rovnici přepíšeme jako $z^3=64i
=64(\cos\frac\pi2+i\sin\frac\pi2)$.
Pro $k=0,1,2$ jsou všechny kořeny
\[
z_k=4\left(\cos\left(\frac\pi6+\frac{2k\pi}3\right)
+i\sin\left(\frac\pi6+\frac{2k\pi}3\right)\right).
\]
Jejich goniometrické argumenty jsou $\frac\pi6$, $\frac{5\pi}6$
a $\frac{3\pi}2$; v algebraickém tvaru dostáváme
\[
z_0=2\sqrt3+2i,\qquad z_1=-2\sqrt3+2i,\qquad z_2=-4i.
\]
\begin{center}
\begin{tikzpicture}[scale=.52]
    \draw[step=1,gray!20,very thin] (-5,-5) grid (5,5);
    \draw[->] (-5.2,0)--(5.3,0) node[right] {$\operatorname{Re}z$};
    \draw[->] (0,-5.2)--(0,5.3) node[above] {$\operatorname{Im}z$};
    \draw[gray] (0,0) circle (4);
    \fill ({2*sqrt(3)},2) circle (2pt) node[above] {$z_0$};
    \fill ({-2*sqrt(3)},2) circle (2pt) node[above] {$z_1$};
    \fill (0,-4) circle (2pt) node[right] {$z_2$};
\end{tikzpicture}
\end{center}
Kontrola: $(2\sqrt3+2i)^3=64i$;
další kořeny získáme násobením třetími kořeny jedničky.
Jde o tři různé kořeny kubické rovnice.
\end{enumerate}

\newpage
\subsection*{Varianta H}
\begin{enumerate}
\item Pro $z=x+yi$ je $z\overline z=x^2+y^2$, tedy
\[
x^2+y^2+x+yi=6+2i.
\]
Odtud $y=2$ a $x^2+x-2=(x-1)(x+2)=0$.
Řešeními jsou $z=1+2i$ a $z=-2+2i$.
Dosazení dává $|1+2i|^2+(1+2i)=5+1+2i=6+2i$ a
$|-2+2i|^2+(-2+2i)=8-2+2i=6+2i$.

\item První podmínka je rovností vzdáleností od bodů $(1,-1)$
a $(3,1)$. Po umocnění obou stran:
\[
(x-1)^2+(y+1)^2=(x-3)^2+(y-1)^2
\quad\Longleftrightarrow\quad x+y=2.
\]
Spolu s $x^2+y^2=10$ dostáváme po dosazení $x=2-y$
\[
(2-y)^2+y^2=10
\quad\Longleftrightarrow\quad (y-3)(y+1)=0.
\]
Řešeními jsou $z=-1+3i$ a $z=3-i$. Oba body splňují
rovnici přímky $x+y=2$ i kružnice $x^2+y^2=10$:
\begin{center}
\begin{tikzpicture}[scale=.58]
    \draw[step=1,gray!20,very thin] (-4,-4) grid (4,4);
    \draw[->] (-4.2,0)--(4.3,0) node[right] {$\operatorname{Re}z$};
    \draw[->] (0,-4.2)--(0,4.3) node[above] {$\operatorname{Im}z$};
    \draw (0,0) circle ({sqrt(10)});
    \draw[thick] (-1.5,3.5)--(3.5,-1.5) node[pos=.6,above right] {$x+y=2$};
    \fill (-1,3) circle (2pt) node[above left] {$-1+3i$};
    \fill (3,-1) circle (2pt) node[below left] {$3-i$};
\end{tikzpicture}
\end{center}

\item Vynásobením rovnice číslem $-i$ získáme
$x^2+3ix+4=0$. Diskriminant je $D=(3i)^2-16=-25$,
odkud
\[
x_{1,2}=\frac{-3i\pm5i}{2},
\qquad x_1=i,\quad x_2=-4i.
\]
Kontrola dosazením do původní rovnice:
$i\cdot i^2-3i+4i=0$ a
$i(-4i)^2-3(-4i)+4i=-16i+12i+4i=0$.

\item Platí $1-\sqrt3i=2(\cos(-\frac\pi3)+i\sin(-\frac\pi3))$ a
$1-i=\sqrt2(\cos(-\frac\pi4)+i\sin(-\frac\pi4))$.
Podíl má modul $\sqrt2$ a argument $-\frac\pi{12}$, proto
\[
\left(\frac{1-\sqrt3i}{1-i}\right)^6
=(\sqrt2)^6\left(\cos(-\frac\pi2)+i\sin(-\frac\pi2)\right)=-8i.
\]
Nezávisle algebraicky: druhá mocnina podílu je $\sqrt3-i$,
jehož třetí mocnina je $-8i$.

\item Rovnice má tvar $z^3=-8i
=8(\cos\frac{3\pi}2+i\sin\frac{3\pi}2)$.
Pro $k=0,1,2$ jsou všechny kořeny
\[
z_k=2\left(\cos\left(\frac\pi2+\frac{2k\pi}3\right)
+i\sin\left(\frac\pi2+\frac{2k\pi}3\right)\right).
\]
Jejich argumenty jsou $\frac\pi2$, $\frac{7\pi}6$
a $\frac{11\pi}6$; v algebraickém tvaru:
\[
z_0=2i,\qquad z_1=-\sqrt3-i,\qquad z_2=\sqrt3-i.
\]
\begin{center}
\begin{tikzpicture}[scale=.65]
    \draw[step=1,gray!20,very thin] (-3,-3) grid (3,3);
    \draw[->] (-3.2,0)--(3.3,0) node[right] {$\operatorname{Re}z$};
    \draw[->] (0,-3.2)--(0,3.3) node[above] {$\operatorname{Im}z$};
    \draw[gray] (0,0) circle (2);
    \fill (0,2) circle (2pt) node[right] {$z_0$};
    \fill ({-sqrt(3)},-1) circle (2pt) node[below] {$z_1$};
    \fill ({sqrt(3)},-1) circle (2pt) node[below] {$z_2$};
\end{tikzpicture}
\end{center}
Kontrola: $(\sqrt3-i)^3=-8i$;
další kořeny získáme násobením třetími kořeny jedničky.
Jde o tři různé kořeny kubické rovnice.
\end{enumerate}
\end{document}
